\section{Introduction}

The machine learning community has long worried about epistemic lock-in: the possibility that reliance on a narrow set of benchmark datasets creates self-reinforcing cycles where early methodological choices become permanently embedded in research practice. This concern appears well-founded---highly cited papers overwhelmingly share benchmark datasets with their references, and concentration on standard evaluation corpora remains high across top venues. Yet despite these surface indicators of rigidity, we find no evidence that the field is actually trapped. Benchmark standards propagate through citation networks, but they do not create the temporal feedback loops that would constitute genuine lock-in.

This distinction matters because it separates legitimate concerns about evaluation methodology from unfounded fears about irrecoverable scientific stagnation. A field where benchmark use correlates with citation patterns differs fundamentally from one where early concentration causally determines future diversity. The former describes co-evolution between community standards and research practice; the latter describes a path-dependent trap. Our contribution is establishing, through quantitative temporal analysis, which of these characterizes machine learning research.

\subsection{The Problem in Three Layers}

At the surface level, benchmark concentration is easily observed. Any researcher attending NeurIPS, ICML, or ICLR notices the same datasets appearing repeatedly: ImageNet for vision, GLUE for language understanding, standard splits on well-known corpora. This repetition raises natural questions about whether the field evaluates methods on sufficiently diverse tasks to generalize findings.

Digging deeper, citation analysis reveals that benchmark choices are not independent across papers. Prior work by \citet{engdahl2024benchmark} documents through ethnographic methods how benchmark standards emerge through community processes, while studies of hyperparameter optimization benchmarks like HPO-B \citep{pineda2021hpob} demonstrate concentrated usage patterns even when alternatives exist. High-impact papers share substantially more evaluation overlap with their references than citation structure alone would predict---a pattern consistent with benchmark standards propagating through intellectual influence.

Yet the deepest question remains unaddressed: does this propagation create a feedback loop? If concentration in year $t$ causally reduces diversity in year $t+1$, and this reduced diversity further increases concentration in year $t+2$, the field would face genuine lock-in requiring intervention. If instead concentration and citation patterns co-evolve without causal feedback---both responding to common factors like dataset availability, community growth, or methodological shifts---then benchmark standardization reflects coordination rather than confinement.

No prior work has conducted the quantitative temporal analysis required to distinguish these hypotheses. This gap motivates our study.

\subsection{Our Approach and Key Insight}

We analyze 12,600 papers from NeurIPS, ICML, and ICLR spanning 2018--2024, using Papers With Code data that provides benchmark usage information with over 80\% coverage. We operationalize concentration through the Herfindahl-Hirschman Index (HHI), computing $\text{HHI} = \sum s_d^2$ where $s_d$ represents the share of papers using dataset $d$.

Our analysis proceeds through six sub-hypotheses testing the mechanism chain from concentration measurement through propagation to temporal dynamics. We first establish that HHI provides a valid concentration measure correlating with simpler top-5 share metrics ($\rho = 0.90$, $p < 0.001$). We then demonstrate that benchmark standards do propagate through citations---papers share substantially higher benchmark overlap with papers they cite (Jaccard similarity 0.318 versus 0.014 baseline, Cohen's $d = 1.93$). Prior concentration predicts individual paper adoption of standard benchmarks (odds ratio $4.4 \times 10^{24}$).

The key finding emerges in our temporal analysis. Despite strong propagation mechanisms, we find no evidence of causal feedback from concentration to subsequent diversity. Panel regression of benchmark entropy on lagged HHI, controlling for venue and year fixed effects, yields a coefficient that is not statistically significant ($\beta = +1.60$, $p = 0.098$). Granger causality tests similarly fail to reject the null. Concentration does not trap the field; benchmark standards spread, but diversity can shift independently.

\subsection{Contributions}

This paper makes three contributions to understanding benchmark dynamics in machine learning research. First, we provide a comprehensive quantitative characterization of dataset concentration across major venues, establishing that while concentration is measurable and substantive, it varies meaningfully across venues and years rather than reflecting uniform rigidity. Second, we document the mechanism by which benchmark standards propagate: citation networks carry evaluation choices, creating correlation in benchmark use between papers and their references far exceeding chance overlap. Third, and most importantly, we test whether this propagation creates self-reinforcing lock-in through panel regression and Granger causality analysis, finding that it does not.
